\documentclass[11pt,a4paper]{article}

\usepackage[a4paper,margin=2.5cm]{geometry}

\usepackage{graphicx}
\usepackage{amsmath}
\usepackage{booktabs}
\usepackage{siunitx}
\usepackage{float}
\usepackage{hyperref}

% Title Section
% -----------------------------------------------------------
\begin{document}
\begin{center}
    \vspace{2cm}

    \Huge{Wave Energy Analysis}\vspace{0.5cm}

    \large{Analysis of Real Ocean Wave Data from the Coast of Portugal Offshore}\vspace{1cm}

    Iuna Dreyer\vspace{0.3cm}

    May 2026\vspace{1cm}
\end{center}


% ------------------------------------------------------------

\section{Introduction}

Ocean wave energy is a promising renewable energy source due to the high energy density of ocean waves and its peak season being winter (which complements solar energy very well).

Portugal possesses favourable Atlantic wave conditions and has been an important region for wave energy research and testing.

This project aims to analyse real ocean wave data obtained from the Copernicus Marine Service to estimate wave energy potential and simulate the response of a simplified wave energy converter model.




\section{Methodology}

\subsection{Wave Data}

The dataset used in this project is the Global Ocean Waves Analysis and Forecast product from the Copernicus Marine Service for the Portuguese offshore region during January, February and March 2026. The variables that were used in this analysis are the significant wave height and the mean wave period, sampled in 3-hour intervals.


\subsection{Dynamic Model}





\section{Results}

To be completed.

% =================================================

\section{Discussion}

To be completed.

% =================================================

\section{Conclusion}

To be completed.

% =================================================

\bibliographystyle{plain}
\bibliography{references}

\end{document}

